% !TEX TS-program = pdflatex
% !TeX program = pdflatex
% !TEX encoding = UTF-8
% !TeX spellcheck = en_US

\documentclass[11pt, a4paper]{article}
%\usepackage{fullpage}
\usepackage[left=1.5cm,right=1.5cm,top=1cm,bottom=2cm]{geometry}
\usepackage[fleqn]{amsmath}
\usepackage{amssymb}
%\usepackage{indentfirst}
\usepackage{fontenc}
\usepackage[utf8]{inputenc}
%\usepackage[arabic,french,english]{babel}
\usepackage{txfonts}
\usepackage[]{graphicx}
\usepackage{multirow}
\usepackage{hyperref}
\usepackage{parskip}
\usepackage{multicol}
\usepackage{wrapfig}
\usepackage{fancyhdr}
\usepackage{hyperref}
\usepackage{tcolorbox}
\usepackage{natbib}
\usepackage{arabtex}

\usepackage{turnstile}%Induction symbole

\renewcommand{\baselinestretch}{1}

\setlength{\parindent}{24pt}

%\usepackage{etoolbox}
%\patchcmd{\thebibliography}{\section*{\refname}}{}{}{}

\fancyhf{}

\lfoot{Abdelkrime Aries}
\cfoot{2023-2024/ESI/2CSSID/NLP}
\rfoot{\textbf{\thepage}}
%
\renewcommand{\headrulewidth}{0pt} % remove lines as well
\renewcommand{\footrulewidth}{1pt}

\newcommand\repeatstr[1]{\leavevmode\xleaders\hbox{#1}\hfill\kern0pt}

\renewcommand{\bibsection}{}
\bibliographystyle{humannat}

\newcommand{\kurl}[1]{{\scriptsize\bfseries\url{#1}}}

\setlength{\parindent}{0pt}


\begin{document}
	

%\selectlanguage {french}
%\pagestyle{empty} 
\pagestyle{fancy}

\noindent
\begin{tabular}{ll}
	\multirow{3}{*}{\includegraphics[width=1.6cm]{../../../extra/logo/esi.nlp.pdf}} & \'Ecole national Supérieure d'Informatique, Algiers\\
	& 2CSSID (2023-2024)\\
	& Natural Language Processing (NLP)
\end{tabular}\\[.25cm]
\noindent\rule{\textwidth}{1pt}\\[-0.25cm]
\begin{center}
	{\LARGE \textbf{Lab02: Tweets' grammaticality judgment}}
	\begin{flushright}
		Abdelkrime Aries
	\end{flushright}
\end{center}\vspace*{-0.25cm}
\noindent\rule{\textwidth}{1pt}

Grammaticality refers to how well a sentence conforms to the rules of a particular language. 
It seeks to judge structure and not meaning. 
In this lab, we want to test Algerian tweets statistically.
In other word, we want to identify the grammatical status of a tweet directly with the probability of its occurrence.
We chose to us N-Grams.

\section{Program description}

Here, different functions are described; either those implemented or not. 
You have to understand how they are implemented to respond on the different questions.

\subsection{Ngram class}

Given an already implemented Bi-grams class (N=2), implement a new one with N predefined.
The class must:
\begin{itemize}
	\item take in consideration the case where N is less than 1 or greater than 6 (must not be accepted);
	\item support both Lidstone and Interpolation smoothing. 
	Let say for both, the parameters are predefined (in case of interpolation, no need to estimate them using another validation dataset);
	\item The parameters must be defined in the constructor;
	\item Interpolation parameters are from the lowest order to the highest;
	\item The smoothing type must be defined while predicting.
\end{itemize}

These are some recommendations concerning \textbf{score} function:
\begin{itemize}
	\item Any division by zero will have a log probability of \textbf{-math.inf}
	\item log probability of 0 is \textbf{-math.inf}
	\item For unigrams, the probability of interpolation is that of Lidstone;
	\item For n-grams $ n>1 $, using interpolation, we consider higher order as 0 when probability is 0 or there is division by zero.
	Example, 0.7 * -inf + 0.3 * 0.02 is 0.06 and not -inf (-inf will be handled as a 0)
	
	\item For easier interpolation, we will not consider $ p(</s>) $.
	In this case, $ p_I(</s>|help\ you) = 0.5 p(</s>|help\ you) + 0.3 p(</s>|you) + 0.2 p(</s>) = 0.5 p(</s>|help\ you) + 0.3 p(</s>|you)$
\end{itemize}

For \textbf{predict} function, anything plus minus infinity is minus infinity

\textbf{\slshape This must be implemented}.

\subsection{Grammaticality class}


\textbf{\slshape Not implemented}.

\section{Questions}

Answer these questions at the start of your code, as comments.
\begin{enumerate}
	\item We want to consider unknown words in general texts, propose a solution (not smoothing, you can use something other than Ngrams).
	We want to take in consideration different variants of Arabizi, propose a solution.
	
	\item If we train a model on the inverse of texts.
	Will we get the same probability as the one in the right direction? Why?
	Will this affect grammaticality judgment? Why?
	
	\item Can we use Viterbi to calculate the probability of a text? Why/How?
	
	\item Describe how can we decide that a text is grammatical or not based on its probability.
	
\end{enumerate}


\section{Evaluation}

\begin{itemize}
	\item Duration: 1h
	\item Grade
	\begin{itemize}
		\item \textbf{Exceptions} (1pt) = $ N < 1 $ (0.5pts) + $ N > 6 $ (0.5pts)
		\item \textbf{fit grade} (6pts) = processing all data (2pts) + all N-grams (2pts) + all less than N-- grams (2pts).
		\item \textbf{score grade} (5pts) = correct and less complex estimation (2pts) + applying smoothing (2pt) + log probability (1pt).
		\item \textbf{predict grade} (3pts) = correctly calculate the log probability (2pts) + correctly taking smoothing in consideration (1pt).
		\item \textbf{questions grade} (4pts) = 1pt for each question.
		\item \textbf{In time grade} (1pt): after the deadline, each late 2 minutes are -0.25. So, 8 minutes then you will get 0.
	\end{itemize}
\end{itemize}

\vfill
\begin{flushright}
	\textit{Ramadan kareem}
\end{flushright}

\end{document}
